% ---------------------------------------------------------------------
% STATUT : RÉDIGÉ (transplanté de exploratory/cascade_qualitative/
%          chapitre_positionnement.tex, préambule retiré).
% RESTE À FAIRE : (1) basculer la liste de références de fin de chapitre
%          dans references.bib et passer aux \citep ; (2) ajouter le
%          positionnement de l'échelle de repli (identification
%          partielle) une fois le chapitre 5 arrêté.
% ---------------------------------------------------------------------
\chapter{État de l'art et positionnement}
\label{ch:positionnement}

La modélisation actuarielle du risque cyber s'est construite en trois vagues, dont ce
mémoire hérite les acquis tout en s'en démarquant sur un point. La première interroge
l'\textbf{assurabilité} même du risque et la rareté des données \citep{BienerElingWirfs2015},
et pose les marges par la théorie des valeurs extrêmes. La deuxième modélise la
\textbf{contagion}, mais entre entités d'un portefeuille, par des modèles de réseau ou
épidémiologiques. La troisième saisit la \textbf{dynamique temporelle} des arrivées par des
processus de Hawkes. Le présent travail emprunte à chacune ses marges empiriques et son
cadre de risque extrême, et emprunte au risque de crédit deux outils structurels que la
littérature cyber n'avait pas mobilisés ensemble : le processus de \textbf{branchement
multitype}, qui décrit la propagation d'un incident entre piliers (chapitre~\ref{chap:cascade}),
et les \textbf{modèles multi-états} de type migration de rating, qui font dépendre le
capital de l'état de conformité (chapitre~\ref{chap:multietats}). Sa démarcation tient en un
mécanisme que ces trois vagues laissent de côté, exposé ci-dessous.

\begin{cle}
\textbf{La contribution en une phrase.} Là où la littérature cyber actuarielle modélise soit
la \emph{co-occurrence} (copules, chocs communs), soit la \emph{contagion entre entités}
(réseaux, épidémiologie), soit la \emph{dynamique temporelle} (Hawkes), ce mémoire introduit
une \textbf{dépendance dirigée et à l'ordre entre les domaines de contrôle réglementaires}
(les piliers DORA) au sein d'une même entité, et la traduit en SCR. Sa signature est que
la criticité dépend de l'\textbf{ordre} de la chaîne et non du seul ensemble des piliers
touchés, ce qu'aucune des approches recensées ne peut reproduire.
\end{cle}

\section{Le mécanisme, et sa signature}

La cascade dirigée pose une matrice $W$ de contagion entre piliers, où $W_{jk}$ mesure la
propension d'une défaillance de $k$ à entraîner celle de $j$. Trois propriétés la
distinguent : elle est \textbf{dirigée} ($W_{jk}\neq W_{kj}$), elle autorise les
\textbf{cycles} (P1$\to$P2$\to\cdots\to$P1), stabilisés par la normalisation de Leontief
($\rho(W)\le g<1$, forme réduite $(I-W)^{-1}$), et surtout elle est \textbf{à l'ordre} : la
criticité d'une chaîne n'est pas une fonction de l'\emph{ensemble} des piliers touchés, elle
dépend de l'\emph{ordre} dans lequel ils le sont. Ce n'est pas une intuition : sur les 26
sous-ensembles de taille au moins deux, \textbf{22 changent de criticité selon l'ordre},
avec un écart atteignant \textbf{6 points sur une échelle de 10} (la chaîne P1$\to$P2 vaut
8, la chaîne P2$\to$P1 vaut 2). Or toute dépendance symétrique et tout score additif sont,
par construction, invariants par permutation : \textbf{aucun ne peut reproduire cela}.

\begin{cle}
\textbf{Remarque de méthode, et ce qui n'est \emph{pas} revendiqué.} Une propriété plus
forte, la \emph{non-transitivité} de la dominance entre piliers, a été testée dans trois
formulations : la transitivité des corrélations qu'impose un modèle à facteur unique, la
positivité de la dépendance symétrisée, et l'existence de cycles de dominance par paire.
\textbf{Aucune ne se vérifie} sur ce modèle. Elle n'est donc pas revendiquée. La
dépendance à l'ordre, elle, se vérifie, et elle exclut exactement les mêmes concurrents.
Signaler l'énoncé que l'on a testé et abandonné vaut mieux que de le laisser à la
sagacité du lecteur.
\end{cle}

\section{Positionnement : pourquoi le mécanisme est irréductible}

\begin{center}\footnotesize
\renewcommand{\arraystretch}{1.35}
\begin{tabular}{@{}p{3.2cm} c c c c p{3.1cm} p{3.0cm}@{}}
\toprule
\textbf{Approche} & \textbf{Dir.} & \textbf{Ordre} & \textbf{Cycles} & \textbf{Propag.} &
\textbf{Donnée requise} & \textbf{Objet modélisé} \\
\midrule
Matrices de risque / AMDEC & \non & \non & \non & \non & dires d'expert & risques listés (proba $\times$ gravité) \\
Copule (Gumbel, gaussienne) & \non & \non & \non & \non & co-mouvement observé & briques de perte \\
Réseau / épidémiologie (SIR) & \oui & \non & \oui & \oui & structure réseau $+$ taux d'infection & entités d'un portefeuille \\
Hawkes (multivarié) & \oui & \non & \oui & \oui & horodatage fin des arrivées & types d'attaque, dans le temps \\
Réseau bayésien & \oui & \non & \non & \oui & tables de proba conditionnelle & variables d'un graphe acyclique \\
\cascell\textbf{Cascade dirigée (ici)} & \cascell\oui & \cascell\oui & \cascell\oui & \cascell\oui & \cascell jugement d'expert (élicité) & \cascell\textbf{domaines de contrôle (piliers) d'une entité} \\
\bottomrule
\end{tabular}
\end{center}

\vspace{2pt}
Point par point :
\begin{itemize}\itemsep2pt
\item \textbf{Matrices de risque / AMDEC} (le standard praticien, RPN $=$ gravité $\times$
      occurrence $\times$ détection) : listent et notent des risques mais n'ont \emph{aucune}
      propagation ni direction. La cascade ajoute exactement la couche qui leur manque.
\item \textbf{Copules} \citep{BohmeKataria2006, HerathHerath2011} : capturent la
      dépendance de queue entre pertes, mais \textbf{symétrique} : elles ne distinguent pas
      $k\to j$ de $j\to k$. Notre benchmark le chiffre : une copule ne peut pas
      exprimer la propagation (interaction moyenne $+0$ par construction contre $+560$ M\euro{}
      pour la cascade), et sa priorité de remédiation suit les \emph{réceptacles}, la cascade
      les \emph{sources}.
\item \textbf{Modèles réseau et épidémiologiques} \citep{HillairetLopez2021,
      HillairetLopezDOultremontSpoorenberg2022, XuHua2017, Baldwin2017} : ils modélisent la
      vraie contagion, mais \textbf{entre
      entités} d'un portefeuille, et \textbf{exigent la structure du réseau} et des taux
      d'infection, données indisponibles ici. Nous \emph{empruntons} à Hillairet et Lopez la
      logique d'états de conformité, mais l'appliquons \textbf{entre piliers} d'une entité,
      pas entre assurés.
\item \textbf{Hawkes} \citep{BessyRolandBoumezouedHillairet2021, BoumezouedCherkaouiHillairet2023} :
      modélisent l'auto-excitation de la \textbf{fréquence dans le temps},
      entre \emph{types d'attaque}, pas la structure inter-piliers. Nous montrons
      (chapitre~\ref{ch:donnees}) que sur la fréquence d'entrée, l'excitation est
      intra-journalière et exogène, et nous ne l'ajoutons pas ; ce constat rejoint leur
      raffinement two-phase de 2023.
\item \textbf{Réseaux bayésiens} : peuvent porter la direction, mais supposent un graphe
      \textbf{acyclique} (pas de cycle de contagion) et des tables de probabilités
      conditionnelles qui se heurtent au \textbf{même mur d'identification} que $W$
      (chapitre~\ref{ch:identifiabilite}). La forme de Leontief, elle, gère les cycles.
\end{itemize}

\section{Revue de littérature, par thème}

\textbf{Dépendance et accumulation.} La modélisation actuarielle du cyber a d'abord traité la
\emph{corrélation} entre risques par copules \citep{BohmeKataria2006, HerathHerath2011}
et par chocs communs. Ces approches capturent l'accumulation (que nos données confirment,
concentration journalière de Gini $0{,}69$) mais restent symétriques.

\textbf{Contagion structurée.} Une seconde vague introduit la propagation : \citet{HillairetLopez2021}
couplent des processus de comptage à un modèle épidémiologique compartimental (type
SIR) pour la contagion dans un portefeuille ; \citet{HillairetLopezDOultremontSpoorenberg2022}
ajoutent une structure de réseau ; \citet{XuHua2017} et \citet{Baldwin2017}
modélisent la contagion entre cibles. Toutes exigent la donnée de réseau ou d'infection.

\textbf{Dynamique temporelle.} La fréquence auto-excitée est traitée par les processus de
Hawkes : \citet{BessyRolandBoumezouedHillairet2021} sur données PRC (dont notre source
dérive), puis \citet{BoumezouedCherkaouiHillairet2023} avec un terme exogène ;
\citet{HillairetReveillacRosenbaum2021} en développent la théorie pour le pricing de
dérivés ; \citet{Peng2016} sur des taux d'attaque horodatés. Notre travail se situe dans
cette lignée mais en \emph{sort} sur la fréquence d'entrée, pour les raisons données au
chapitre~\ref{ch:donnees}.

\textbf{Fréquence, sévérité, queue.} Sur les marges, \citet{Edwards2016} ajustent des lois
discrètes (binomiale négative) au nombre journalier de brèches ; \citet{LuZhangZhu2024}
séparent fréquence et nombre de brèches par événement (structure de paquet qui soutient notre
Poisson composé) ; \citet{FarkasLopezThomas2021} tarifent par arbres de régression GPD. Nous
en retenons la théorie des valeurs extrêmes (indice de queue $\xi$) et la structure de paquet.

\textbf{Cadre structurel.} Le squelette quantitatif emprunte au modèle à facteur de
\citet{Vasicek2002} (risque de crédit) l'idée d'un stress latent à composante systémique, et
à l'analyse entrées-sorties de \citet{Leontief1936} la normalisation qui garantit la
stabilité d'une contagion
\emph{avec cycles}. La nouveauté est de faire porter cette mécanique par une matrice
\textbf{dirigée} entre domaines de contrôle réglementaires.

\section{Où se situe ce mémoire}

Le mémoire \textbf{emprunte} ce qui est établi et ancré sur données (la queue de sévérité par
EVT, la structure de paquet de la fréquence, la logique d'états de conformité de Hillairet et
Lopez) et \textbf{ajoute} deux choses absentes de la littérature : (i)~une dépendance
\textbf{dirigée et à l'ordre} entre les piliers DORA, traduite en
capital ; (ii)~une \textbf{frontière d'identifiabilité} explicite, muée en spécification de
reporting. Aucune des références recensées ne modélise une dépendance dirigée entre domaines
de contrôle au sein d'une entité, ni ne la relie au SCR.

\begin{cle}
\textbf{La revendication de nouveauté, défendable.} \og Je n'invente pas une nouvelle loi de
sévérité ni un nouveau processus de fréquence : sur ces marges, je m'appuie sur la
littérature établie et je la calibre. Ma contribution est le \textbf{mécanisme de dépendance}
lui-même, dirigé et à l'ordre entre piliers réglementaires, que ni les copules (symétriques),
ni les modèles de réseau (qui exigent une donnée absente), ni les processus de Hawkes (qui
vivent dans le temps) ne capturent, et que je traduis en capital tout en démontrant sa limite
d'identification. \fg{} C'est une nouveauté \emph{maîtrisée} : neuve sur le mécanisme, sobre
sur le reste, honnête sur les limites.
\end{cle}

% Les références de ce chapitre sont dans references.bib et rendues par la bibliographie
% de fin. Le \nocite garantit l'apparition de celles qui ne seraient pas déjà \citep dans
% le corps (filet de sécurité).
\nocite{BohmeKataria2006, HerathHerath2011, HillairetLopez2021,
  HillairetLopezDOultremontSpoorenberg2022, XuHua2017, Baldwin2017,
  BessyRolandBoumezouedHillairet2021, BoumezouedCherkaouiHillairet2023,
  HillairetReveillacRosenbaum2021, Peng2016, Edwards2016, LuZhangZhu2024,
  FarkasLopezThomas2021, Vasicek2002, Leontief1936}

% =====================================================================
